\documentclass[11pt,a4paper,fontset=none]{ctexart}

\usepackage[a4paper,top=22mm,bottom=22mm,left=23mm,right=23mm]{geometry}
\usepackage{fontspec}
\usepackage{amsmath,amssymb}
\usepackage{booktabs,longtable,tabularx,array}
\usepackage{xcolor}
\usepackage{enumitem}
\usepackage{listings}
\usepackage{fancyhdr}
\usepackage{hyperref}

\setmainfont{Arial}
\setsansfont{Arial}
\setmonofont{Consolas}
\setCJKmainfont{Microsoft YaHei}
\setCJKsansfont{Microsoft YaHei}
\setCJKmonofont{Microsoft YaHei}

\definecolor{SustechOrange}{HTML}{F36C21}
\definecolor{SustechGreen}{HTML}{006D68}
\definecolor{Ink}{HTML}{1F2933}
\definecolor{Muted}{HTML}{5F6B76}
\definecolor{Pale}{HTML}{F4F7F8}
\definecolor{Warn}{HTML}{FFF4EA}

\hypersetup{
  colorlinks=true,
  linkcolor=SustechGreen,
  urlcolor=SustechGreen,
  pdftitle={MH6手掌闭环求解器程序说明},
  pdfauthor={Jiarui Su},
  pdfsubject={解析闭环求解、工作空间映射与安全接口说明}
}

\ctexset{
  section={format=\Large\bfseries\color{SustechGreen},beforeskip=1.7ex,afterskip=0.7ex},
  subsection={format=\large\bfseries\color{Ink},beforeskip=1.2ex,afterskip=0.4ex}
}

\pagestyle{fancy}
\setlength{\headheight}{13pt}
\fancyhf{}
\fancyhead[L]{\small\color{Muted}MH6 手掌闭环求解器}
\fancyhead[R]{\small\color{Muted}改进版程序说明}
\fancyfoot[C]{\small\color{Muted}\thepage}
\renewcommand{\headrulewidth}{0.3pt}

\setlist[itemize]{leftmargin=1.6em,itemsep=0.25em,topsep=0.35em}
\setlist[enumerate]{leftmargin=1.8em,itemsep=0.3em,topsep=0.35em}
\renewcommand{\arraystretch}{1.22}

\lstdefinestyle{pythonstyle}{
  language=Python,
  basicstyle=\ttfamily\small,
  keywordstyle=\color{SustechGreen}\bfseries,
  commentstyle=\color{Muted},
  stringstyle=\color{SustechOrange},
  backgroundcolor=\color{Pale},
  frame=single,
  rulecolor=\color{black!15},
  breaklines=true,
  showstringspaces=false,
  columns=fullflexible,
  xleftmargin=0.5em,
  framexleftmargin=0.4em
}

\newenvironment{summarybox}{%
  \par\smallskip\noindent\begin{minipage}{\linewidth}%
  \color{SustechGreen}\rule{\linewidth}{0.8pt}\par\smallskip\color{Ink}%
}{%
  \par\smallskip\color{SustechGreen}\rule{\linewidth}{0.8pt}%
  \end{minipage}\par\smallskip%
}
\newenvironment{warningbox}{%
  \par\smallskip\noindent\begin{minipage}{\linewidth}%
  \color{SustechOrange}\rule{\linewidth}{0.8pt}\par\smallskip\color{Ink}%
}{%
  \par\smallskip\color{SustechOrange}\rule{\linewidth}{0.8pt}%
  \end{minipage}\par\smallskip%
}

\begin{document}

\begin{titlepage}
  \centering
  \vspace*{28mm}
  {\color{SustechOrange}\rule{\textwidth}{3pt}}\\[15mm]
  {\Huge\bfseries\color{Ink} MH6 手掌闭环求解器\\[4mm]程序说明}\par
  \vspace{8mm}
  {\Large\color{SustechGreen}解析闭环求解、工作空间映射与安全接口}\par
  \vfill
  \begin{summarybox}
    \centering
    \textbf{版本：}改进版 v1.1-safety\quad
    \textbf{日期：}2026-09-20\\[2mm]
    \textbf{运行环境：}Python 3.12 / NumPy 2.3.5
  \end{summarybox}
  \vfill
  {\large Jiarui Su}\par
  \vspace{8mm}
  {\color{SustechOrange}\rule{\textwidth}{3pt}}
\end{titlepage}

\tableofcontents
\newpage

\section{项目定位}

本程序针对 MH6 手掌球面六杆闭环机构。给定三个已知输入角

\[
q_{\mathrm{in}}=[\alpha_2,\alpha_3,\theta_1]^T,
\]

解析求解三个闭环从属角

\[
q_{\mathrm{dep}}=[\theta_2,\theta_3,\alpha_1]^T,
\]

并根据三路独立标定关系输出电机控制值。程序还提供归一化输入映射、闭环可行性诊断、无解轨迹投影和双分支连续选择。

\begin{summarybox}
本程序是\textbf{闭环补全求解器}，不是根据末端笛卡尔位姿求解关节角的通用逆运动学，也不包含手指雅可比、动力学或实机反馈控制。
\end{summarybox}

\subsection{主要能力}

\begin{itemize}
  \item 解析求解球面六杆闭环的两组装配分支；
  \item 同时检查旋转闭环和平移闭环误差；
  \item 区分输入超限与刚性闭环无解；
  \item 提供三种 $[0,1]^3$ 归一化映射模式；
  \item 将无解角度局部投影到可行域内部，并显式报告角度偏差；
  \item 根据上一帧电机值选择变化较小的解析分支；
  \item 支持 API 电机顺序和历史时序 CSV 顺序；
  \item 提供单元测试、网格扫描和轨迹实验脚本。
\end{itemize}

\section{数学模型与解析求解}

\subsection{坐标链和闭环条件}

六段局部旋转写为

\[
\begin{aligned}
R_1&=R_y(80^\circ)R_z(\alpha_2), &
R_2&=R_y(-80^\circ)R_z(\theta_1),\\
R_3&=R_y(80^\circ)R_z(\theta_2), &
R_4&=R_y(35^\circ)R_z(\theta_3),\\
R_5&=R_y(20^\circ)R_z(\alpha_3), &
R_6&=R_y(225^\circ)R_z(\alpha_1).
\end{aligned}
\]

刚性闭环要求

\[
R_1R_2R_3R_4R_5R_6=I.
\]

令 $P=(R_1R_2)^T$，并取 $C=R_y(-80^\circ)P$。由于末端的 $R_z(\alpha_1)$ 不改变旋转矩阵第三列，可以先从第三列约束中消去 $\alpha_1$。

\subsection{求解顺序}

定义

\[
v=R_y(20^\circ)R_z(\alpha_3)R_y(225^\circ)e_3,
\qquad e_3=[0,0,1]^T.
\]

由第三列的 $z$ 分量得到

\[
\cos(\theta_3+\phi)=
\frac{\cos35^\circ\,v_z-C_{33}}
{\sin35^\circ\sqrt{v_x^2+v_y^2}},
\qquad
\phi=\operatorname{atan2}(v_y,v_x).
\]

当右侧绝对值不大于 1 时，通常得到两组 $\theta_3$。随后依次计算

\[
\theta_3\longrightarrow\theta_2\longrightarrow\alpha_1.
\]

每组候选解都会重新计算完整旋转闭环和平移闭环误差。当前程序返回

\[
(\theta_2,\theta_3,\alpha_1,e_R,e_p),
\]

其中 $e_R$ 为旋转矩阵最大元素误差，$e_p$ 为闭环平移残差范数。

\subsection{输入限位}

\begin{center}
\begin{tabular}{lcc}
\toprule
输入 & 下限 & 上限\\
\midrule
$\alpha_2$ & $-31.1^\circ$ & $90.8^\circ$\\
$\alpha_3$ & $-180^\circ$ & $59^\circ$\\
$\theta_1$ & $-23.7^\circ$ & $8.6^\circ$\\
\bottomrule
\end{tabular}
\end{center}

角度分别位于限位内并不保证组合后能够闭合。解析判据为

\[
|c_{\mathrm{closure}}|\leq1.
\]

\section{输入映射与工作空间}

\subsection{三种归一化模式}

\noindent\begin{tabularx}{\textwidth}{>{\ttfamily\raggedright\arraybackslash}p{0.30\textwidth}X>{\raggedleft\arraybackslash}p{0.14\textwidth}}
\toprule
模式 & 含义 & $21^3$ 网格有解率\\
\midrule
workspace\_conditional & 默认。根据 $\alpha_2,\theta_1$ 动态计算可闭合的 $\alpha_3$ 区间 & 100.0\%\\
workspace\_independent & 三个输入分别映射到固定保守区间 & 92.9\%\\
motor\_range & 兼容旧版，独立覆盖完整标定行程 & 20.9\%\\
\bottomrule
\end{tabularx}

默认条件映射使用

\[
\alpha_2=5^\circ+80^\circ u_1,
\qquad
\theta_1=-2^\circ-21.7^\circ u_3,
\]

再根据这两个角解析计算满足

\[
|c_{\mathrm{closure}}|\leq0.999
\]

的最宽 $\alpha_3$ 区间，并将 $u_2$ 映射到该区间。

\begin{warningbox}
条件映射的 100\% 表示\textbf{重新定义后的归一化输入域全部可解}，不表示完整电机行程的笛卡尔积全部可达。$u_2$ 的物理含义依赖 $u_1$ 和 $u_3$，因此不再是固定的第二路电机行程百分比。
\end{warningbox}

\section{程序接口}

\subsection{直接角度求解}

\begin{lstlisting}[style=pythonstyle]
from solve_from_arpha2_arpha3_theta1 import MH6PalmSolver

solver = MH6PalmSolver()

# 输入单位均为 deg
solutions = solver.solve_remaining(47, -80, -20)
# 每组: theta2, theta3, alpha1, rotation_error, translation_error

motor_solutions = solver.solve_motor(47, -80, -20)
# API 顺序: [motor_alpha1, motor_alpha2, motor_alpha3]
\end{lstlisting}

\subsection{归一化输入}

\begin{lstlisting}[style=pythonstyle]
motors = solver.solve_motor_from_normalized(
    0.641, 0.582, 0.885,
    mode="workspace_conditional",
)

# 恢复旧版完整电机行程映射
legacy = solver.solve_motor_from_normalized(
    0.641, 0.582, 0.885,
    mode="motor_range",
)
\end{lstlisting}

\newpage
\subsection{推荐安全入口}

改进版默认不自动修改无解动作：

\begin{lstlisting}[style=pythonstyle]
result = solver.solve_motor_safe(
    alpha2,
    alpha3,
    theta1,
    previous_motor=last_motor,
    project_invalid=False,
)

if not result["success"]:
    raise RuntimeError(result["error"])

command = result["selected"]
\end{lstlisting}

只有业务明确允许修改轨迹时，才显式开启投影，并给出三轴偏差上限：

\begin{lstlisting}[style=pythonstyle]
result = solver.solve_motor_safe(
    alpha2,
    alpha3,
    theta1,
    previous_motor=last_motor,
    project_invalid=True,
    enforce_margin=True,
    target_abs_value=0.999,
    max_projection_delta_deg=[2.0, 5.0, 1.0],
    motor_bounds=(0.0, 1000.0),
)

if not result["success"]:
    raise RuntimeError(result["error"])
\end{lstlisting}

上例的偏差上限仅说明参数顺序，不是经过实机认证的推荐值。

\subsection{关键返回字段}

\begin{longtable}{>{\raggedright\arraybackslash\ttfamily\small}p{0.36\textwidth}p{0.55\textwidth}}
\toprule
字段 & 含义\\
\midrule
\endhead
success & 是否最终得到可接受的电机解\\
requested\_angles & 调用方原始输入角\\
canonical\_requested\_angles & 周期角规范到 $[-180^\circ,180^\circ)$ 后的输入\\
used\_angles & 实际用于求解的角度\\
projected & 是否执行了工作空间投影\\
projection.angle\_delta\_deg & 三轴投影修改量，顺序为 $[\alpha_2,\alpha_3,\theta_1]$\\
projection\_within\_limits & 投影偏差是否通过调用方阈值\\
solutions & 通过电机范围检查的全部分支\\
rejected\_motor\_solutions & 因控制值越界而被拒绝的分支\\
selected & 根据上一帧连续性选择的候选命令\\
diagnostic & 原输入的限位和闭环诊断\\
error & 失败原因\\
\bottomrule
\end{longtable}

\section{电机标定与顺序}

程序当前使用三组独立线性标定：

\[
\begin{aligned}
m_1 &= 500 + \frac{99}{23.6}\alpha_1,\\
m_2 &= 500 - \frac{380}{90.8}\alpha_2,\\
m_3 &= 247 - \frac{753}{180}\alpha_3.
\end{aligned}
\]

这些是项目控制值，不应自行解释为电机角度、脉冲数或编码器计数。近似系数 $0.239$ 的正确方向是

\[
\Delta\text{angle}\approx0.239\,\Delta\text{control},
\qquad
\Delta\text{control}\approx\frac{\Delta\text{angle}}{0.239}.
\]

三轴仍应以各自的零位和标定斜率为准。

默认 API 输出顺序为

\[
[m_{\alpha_1},m_{\alpha_2},m_{\alpha_3}],
\]

历史时序 CSV 使用

\[
[m_{\alpha_3},m_{\alpha_2},m_{\alpha_1}].
\]

调用时可用 \texttt{motor\_order="timeseries"} 统一重排输出和上一帧参考向量。

\section{改进版变更}

\begin{enumerate}
  \item \textbf{修复周期角标定不一致。}例如 $\alpha_3=280^\circ$ 与 $-80^\circ$ 在旋转上等价。旧代码可能允许前者通过闭环检查，却仍以 $280^\circ$ 计算电机值。改进版在求解、诊断和标定前统一规范周期角。
  \item \textbf{禁止默认静默投影。}\texttt{solve\_motor\_safe()} 的 \texttt{project\_invalid} 默认值由 \texttt{True} 改为 \texttt{False}。
  \item \textbf{增加投影偏差硬阈值。}新增 \texttt{max\_projection\_delta\_deg}；任一轴超限时返回失败，不生成候选命令。
  \item \textbf{增加电机控制值范围检查。}默认仅保留三路控制值均在 $[0,1000]$ 内的分支，并报告被拒绝分支。
  \item \textbf{增加非有限值检查。}电机解重排前必须确认三个数均为有限值。
  \item \textbf{修正标定说明。}明确 $0.239$ 是近似的角度变化/控制值变化，而不是绝对控制值的直接乘法公式。
  \item \textbf{补充回归测试。}覆盖安全默认值、周期角等价性、投影偏差拒绝和电机范围过滤。
\end{enumerate}

\section{验证结果}

\subsection{自动化测试}

\begin{center}
\noindent\begin{tabularx}{\textwidth}{Xr}
\toprule
验证项目 & 结果\\
\midrule
单元测试 & 14/14 通过\\
Python 语法编译 & 通过\\
条件映射 $21^3$ 网格 & 9261/9261 有解\\
网格解析分支数 & 18522\\
最大旋转闭环误差 & $9.15\times10^{-16}$\\
最大平移闭环误差 & $6.84\times10^{-6}$ mm\\
随机条件映射样本 & 5000/5000 有解\\
随机生产求解样本 & 1000/1000 均得到两组解\\
\bottomrule
\end{tabularx}
\end{center}

\subsection{302 帧轨迹实验}

\begin{itemize}
  \item 原轨迹无解帧：133/302；
  \item 为处理无解点和近边界点，共投影 135 帧；
  \item 投影后：302/302 有解；
  \item $\alpha_3$ 最大修改量：$25.54^\circ$；
  \item 连续分支选择的电机帧间变化中位数：1.84；
  \item 始终选第一分支时的中位数：12.20。
\end{itemize}

\begin{warningbox}
“投影后全部有解”只证明数学闭环可实现，不证明投影后的动作仍符合原抓握意图。$25.54^\circ$ 的最大角度修改说明实机调用必须设置项目审批过的偏差阈值。
\end{warningbox}

\section{已知限制}

当前程序没有实现以下功能：

\begin{itemize}
  \item 碰撞和机械干涉检测；
  \item 电机速度、加速度、力矩和温度限制；
  \item 传感器反馈、实时通信和闭环控制；
  \item 全局最近可行点投影；
  \item 动力学、接触力和结构柔顺性；
  \item 对实机安全距离的认证。
\end{itemize}

当前投影算法沿闭合判据的局部梯度迭代，不保证得到全局最近点。连续分支选择只比较三路控制值的未加权欧氏距离，也不等价于完整的速度规划。

\section{集成与实机验证清单}

\begin{enumerate}
  \item 明确每个上层调用使用 \texttt{workspace\_conditional} 还是兼容模式；
  \item 对所有外部角度先运行 \texttt{diagnose\_input()}；
  \item 默认关闭投影，确需投影时设置三轴偏差阈值；
  \item 检查 \texttt{projected}、\texttt{used\_angles} 和 \texttt{error}，禁止只读取 \texttt{selected}；
  \item 固化电机顺序、零位、标定斜率和控制值范围；
  \item 增加速度、加速度、碰撞、力矩和急停保护；
  \item 建立旧版本回滚点；
  \item 先进行低速、无负载、有人监护的实机验证；
  \item 实机验证通过后再确定 \texttt{target\_abs\_value} 和投影偏差阈值。
\end{enumerate}

\section{文件说明与复验命令}

\small
\noindent\begin{tabularx}{\textwidth}{>{\ttfamily\raggedright\arraybackslash}p{0.45\textwidth}X}
\toprule
文件 & 用途\\
\midrule
solve\_from\_arpha2\_arpha3\_theta1.py & 正式解析求解器\\
test\_mh6\_palm\_solver.py & 单元和回归测试\\
audit\_conditional\_grid.py & 默认条件映射的全网格验收\\
workspace\_mapping\_experiment.py & 解析区间、随机覆盖率和性能实验\\
project\_trajectory\_experiment.py & 302 帧轨迹投影和连续选解实验\\
\_scan\_full.py & 三种归一化模式的覆盖率扫描\\
README.md & 文本版使用说明\\
\bottomrule
\end{tabularx}
\normalsize

\begin{lstlisting}[style=pythonstyle]
python -B -m unittest -v test_mh6_palm_solver.py
python -B audit_conditional_grid.py --step 0.05
python -B workspace_mapping_experiment.py
python -B project_trajectory_experiment.py
\end{lstlisting}

\begin{summarybox}
\textbf{交付状态：}本版本已经完成离线解析验证、网格扫描和回归测试，可用于代码评审与集成准备；在补齐实机限速、碰撞、力矩、急停和监护验证前，不应作为已放行的实机控制器。
\end{summarybox}

\end{document}
